% Section 5.9, translated from sv/chapters/05/exercises.tex.

\section{Exercises}
\label{c5:sec:uppgifter}

\begin{aside}
\textbf{How to work with the exercises.} Parts~1 and~2 are done during the lesson: first on your
own, a few minutes per exercise, then together. Part~3 is extra exercises for further practice
after the lesson. Most of them can be done in your head or with pencil and paper.

\facitnotis{L05}
\end{aside}

% ------------------------------------------------------------------------------------------
\uppgiftsdel{Part 1}{Review exercises}

\uppgift{1.1}{System of equations: voltage and current}
Solve the following system of equations (currents in amperes):
\[
  \begin{cases}
    2I_1 + 3I_2 = 13 \\
    4I_1 - I_2 = 5
  \end{cases}
\]

\uppgift{1.2}{Circuit analysis with KVL and KCL}
A circuit has three resistors in a network. KVL and KCL give the following system for $U_A$ and
$U_B$ (node voltages in volts):
\[
  \begin{cases}
    3U_A - U_B = 12 \\
    -U_A + 4U_B = 6
  \end{cases}
\]
Solve the system and give $U_A$ and $U_B$.

% ------------------------------------------------------------------------------------------
\uppgiftsdel{Part 2}{New material}

\uppgift{2.1}{Laws of exponents}
Simplify with the laws of exponents. Leave the answer with positive exponents.
\begin{delar}(5)
  \task $x^3 \cdot x^5$
  \task $\dfrac{R^6}{R^2}$
  \task $(I^2)^3$
  \task $U^{-2} \cdot U^5$
  \task $\left(\dfrac{2}{R}\right)^3$
\end{delar}

\uppgift{2.2}{Scientific notation for component values}
Write the following values in scientific notation, $a \times 10^n$ with $1 \leq a < 10$.
\begin{delar}(4)
  \task $47\,000\,\Omega$
  \task $0.000\,022\,\text{F}$
  \task $3\,300\,000\,\text{Hz}$
  \task $0.000\,000\,015\,\text{A}$
\end{delar}

\uppgift{2.3}{Calculating power}
Calculate the power $P$ dissipated in a resistor using $P = \dfrac{U^2}{R}$ and $P = RI^2$.
\begin{parts}
  \item $U = 5\,\text{V}$ and $R = 100\,\Omega$. Calculate $P$ in mW.
  \item $I = 40\,\text{mA}$ and $R = 50\,\Omega$. Calculate $P$ in mW.
\end{parts}

\uppgift{2.4}{RMS value of a sine voltage}
The relationship between the amplitude $\abs{U}$ and the RMS value $U_{\text{RMS}}$ of a sine
voltage is
\[
  U_{\text{RMS}} = \frac{\abs{U}}{\sqrt{2}}.
\]
\begin{parts}
  \item A sine voltage has the amplitude $\abs{U} = 325\,\text{V}$ (mains voltage). Calculate
    $U_{\text{RMS}}$.
  \item A measured RMS voltage is $U_{\text{RMS}} = 12\,\text{V}$. Calculate the amplitude
    $\abs{U}$.
\end{parts}

% ------------------------------------------------------------------------------------------
\uppgiftsdel{Part 3}{Extra exercises}
The exercises below are extra practice, best done on your own after the lesson.

\uppgift{3.1}{Evaluating powers}
Calculate \textbf{without} a calculator.
\begin{delar}(6)
  \task $3^4$
  \task $(-2)^5$
  \task $10^{-2}$
  \task $5^0$
  \task $2^{-3}$
  \task $0.1^3$
\end{delar}

\uppgift{3.2}{Laws of exponents}
Simplify. The answer should have positive exponents.
\begin{delar}(6)
  \task $a^7 \cdot a^{-3}$
  \task $\dfrac{U^5}{U^8}$
  \task $(2R^3)^2$
  \task $\dfrac{(I^2)^4}{I^3}$
  \task $\left(\dfrac{R^2}{3}\right)^3$
  \task $\dfrac{6x^5}{2x^2}$
\end{delar}

\uppgift{3.3}{Powers of ten}
Write as a single power of ten.
\begin{delar}(5)
  \task $10^4 \cdot 10^{-7}$
  \task $\dfrac{10^{-2}}{10^{-5}}$
  \task $\left(10^3\right)^{-2}$
  \task $\dfrac{10^6 \cdot 10^{-2}}{10^{-3}}$
  \task $\sqrt{10^{-6}}$
\end{delar}

\uppgift{3.4}{Roots}
\begin{parts}
  \item Calculate $\sqrt{81}$.
  \item Calculate $\sqrt{0.25}$.
  \item Calculate $\sqrt[3]{27}$.
  \item Calculate $\sqrt{16 \times 10^{-4}}$.
  \item Simplify $\sqrt{50}$ to exact form.
  \item Show with $a = 9$ and $b = 16$ that $\sqrt{a + b} \neq \sqrt{a} + \sqrt{b}$.
\end{parts}

\uppgift{3.5}{Fractional exponents}
Calculate.
\begin{delar}(5)
  \task $16^{1/2}$
  \task $27^{1/3}$
  \task $9^{3/2}$
  \task $8^{-1/3}$
  \task $\left(\dfrac{1}{4}\right)^{-1/2}$
\end{delar}

\uppgift{3.6}{From prefix to base unit}
Rewrite in the base unit, in scientific notation.
\begin{delar}(6)
  \task $2.2\,\text{k}\Omega$
  \task $470\,\mu\text{F}$
  \task $15\,\text{pF}$
  \task $250\,\text{mA}$
  \task $1.8\,\text{MHz}$
  \task $6.8\,\text{nH}$
\end{delar}

\uppgift{3.7}{From scientific notation to prefix}
Give the value with the prefix that gives a number between $1$ and $1000$.
\begin{delar}(3)
  \task $3.3 \times 10^{-6}\,\text{F}$
  \task $4.7 \times 10^{5}\,\Omega$
  \task $2.5 \times 10^{-4}\,\text{A}$
  \task $8 \times 10^{7}\,\text{Hz}$
  \task $1.2 \times 10^{-11}\,\text{F}$
\end{delar}

\uppgift{3.8}{Calculating in scientific notation}
Calculate and give the answer in scientific notation.
\begin{delar}(3)
  \task $\left(4 \times 10^5\right)\left(2 \times 10^{-8}\right)$
  \task $\dfrac{9 \times 10^{-3}}{3 \times 10^{2}}$
  \task $\left(5 \times 10^3\right)^2$
  \task $2.5 \times 10^{-3} + 7.5 \times 10^{-4}$
  \task $\dfrac{\left(6 \times 10^{-2}\right)^2}{4 \times 10^{-5}}$
\end{delar}

\uppgift{3.9}{Significant figures and rounding}
\begin{parts}
  \item How many significant figures does $4.70 \times 10^3$ have?
  \item Round $0.024\,68$ to two significant figures.
  \item Round $12\,345\,\Omega$ to three significant figures and write the answer in scientific
    notation.
  \item A voltage is measured as $U = 5.0\,\text{V}$ and a resistance as $R = 100\,\Omega$ (three
    significant figures). With how many significant figures should the current
    $I = \dfrac{U}{R}$ be given? Calculate the current.
\end{parts}

\uppgift{3.10}{Ohm's law with powers of ten}
Calculate and give the answer with a suitable prefix.
\begin{parts}
  \item $U = 3.3\,\text{V}$ and $R = 220\,\Omega$. Calculate $I$.
  \item $I = 250\,\mu\text{A}$ and $R = 47\,\text{k}\Omega$. Calculate $U$.
  \item $U = 12\,\text{V}$ and $I = 4\,\text{mA}$. Calculate $R$.
\end{parts}

\uppgift{3.11}{Power with powers of ten}
\begin{parts}
  \item $R = 1.5\,\text{k}\Omega$ and $I = 20\,\text{mA}$. Calculate $P = RI^2$.
  \item $U = 5\,\text{V}$ and $R = 2.2\,\text{k}\Omega$. Calculate $P = \dfrac{U^2}{R}$ in mW.
  \item A resistor is rated $470\,\Omega$ and $0.25\,\text{W}$. Calculate the largest permitted
    current.
  \item Calculate the largest permitted voltage across the resistor in c).
\end{parts}

\uppgift{3.12}{Time constant with prefixes}
The time constant of an RC circuit is $\tau = RC$.
\begin{parts}
  \item $R = 10\,\text{k}\Omega$ and $C = 4.7\,\mu\text{F}$. Calculate $\tau$ in ms.
  \item $\tau = 1\,\text{ms}$ and $C = 100\,\text{nF}$. Calculate $R$.
  \item $R = 2.2\,\text{M}\Omega$ and $\tau = 1.1\,\text{s}$. Calculate $C$ and give the answer in
    nF.
\end{parts}

\uppgift{3.13}{Peak value, RMS value and power}
For a sine voltage, $U_{\text{RMS}} = \dfrac{\abs{U}}{\sqrt{2}}$.
\begin{parts}
  \item Calculate $U_{\text{RMS}}$ when the amplitude is $\abs{U} = 10\,\text{V}$.
  \item Calculate the amplitude $\abs{U}$ when $U_{\text{RMS}} = 230\,\text{V}$.
  \item The power in a resistor is calculated from the RMS value, $P = \dfrac{U_{\text{RMS}}^2}{R}$.
    Show that this can be written $P = \dfrac{\abs{U}^2}{2R}$.
  \item Calculate $P$ when $\abs{U} = 12\,\text{V}$ and $R = 8\,\Omega$.
\end{parts}

\uppgift{3.14}{Find the error}
Each line contains a common mistake with powers. Explain the mistake and write the correct answer.
\begin{delar}(3)
  \task $2^3 \cdot 2^4 = 2^{12}$
  \task $\left(a^2\right)^3 = a^5$
  \task $\dfrac{10^6}{10^2} = 10^3$
  \task $(2a)^3 = 2a^3$
  \task $\sqrt{9 + 16} = 3 + 4$
  \task $a^{-2} = -a^2$
\end{delar}
